\documentclass[a4paper]{article}

%------------------------------------------------
% الحزم
%------------------------------------------------

\usepackage{enumitem}
\usepackage[margin=2.5cm]{geometry}
\usepackage{amsmath,amsfonts,amssymb}
\usepackage[charter]{mathdesign}
\usepackage{graphicx}
\usepackage{array}
\usepackage{xcolor}
\usepackage{needspace}
\usepackage{fancyhdr}
\usepackage{polyglossia}

%------------------------------------------------
% اللغات والخطوط
%------------------------------------------------

\setdefaultlanguage[
    calendar=gregorian,
    numerals=maghrib
]{arabic}

\setotherlanguage{english}

% إذا لم تكن الخطوط مثبتة في الجهاز، تُؤخذ من المجلد ../../fonts/
\IfFontExistsTF{Amiri}{%
    \newfontfamily\arabicfont[Script=Arabic, Scale=1.2]{Amiri}
}{%
    \newfontfamily\arabicfont[
        Script=Arabic, Scale=1.2,
        Path=../../fonts/, Extension=.ttf,
        UprightFont=*-Regular, BoldFont=*-Bold,
        ItalicFont=*-Italic, BoldItalicFont=*-BoldItalic
    ]{Amiri}
}

\IfFontExistsTF{AlBattar}{%
    \newfontfamily\arabicfontsf[Script=Arabic, Scale=1.3]{AlBattar}
    \newfontfamily\khatt[Script=Arabic, Scale=1.3]{AlBattar}
}{%
    \newfontfamily\arabicfontsf[
        Script=Arabic, Scale=1.3, Path=../../fonts/,
        BoldFont=ae_AlBattar.ttf
    ]{ae_AlBattar.ttf}
    \newfontfamily\khatt[
        Script=Arabic, Scale=1.3, Path=../../fonts/,
        BoldFont=ae_AlBattar.ttf
    ]{ae_AlBattar.ttf}
}

%------------------------------------------------
% نسخة الطالب / نسخة التصحيح
% لإنتاج نسخة التصحيح: عرّف \withanswers قبل \input لهذا الملف
%------------------------------------------------

\newif\ifcorrection
\ifdefined\withanswers \correctiontrue \else \correctionfalse \fi

\newcommand{\emptybox}{$\square$}
\ifcorrection
    \newcommand{\correctbox}{\textcolor{green!50!black}{$\boxtimes$}}
    \newcommand{\correct}[1]{\textcolor{green!40!black}{\textbf{#1}}}
    \newcommand{\explanation}[1]{%
        \par\vspace{3pt}
        \fcolorbox{gray!60}{gray!8}{\parbox{\dimexpr\linewidth-2\fboxsep-2\fboxrule}{%
        \small\textbf{التعليل:} #1}}}
\else
    \newcommand{\correctbox}{$\square$}
    \newcommand{\correct}[1]{#1}
    \newcommand{\explanation}[1]{}
\fi

\newcommand{\pts}[1]{\textbf{(#1 ن)}}

%------------------------------------------------
% إعداد الصفحة
%------------------------------------------------

\renewcommand{\footrulewidth}{1pt}
\renewcommand{\headrulewidth}{1pt}

\rhead{\khatt المدرسة العليا للأساتذة ببوزريعة}
\lhead{\khatt ملحقة بومرداس}

\cfoot{\thepage}
\lfoot{\khatt الأستاذ: د. بن عبد الرزاق يوسف}

\parindent 0pt
\setlength{\headsep}{15pt}
\arraycolsep=1.4pt

\pagestyle{fancy}

\setlist[enumerate,1]{label=\textbf{السؤال \arabic*:}, leftmargin=*, itemsep=12pt, align=left}

%================================================
% بداية الوثيقة
%================================================

\begin{document}

\begin{center}

{\Large\khatt\textbf{الامتحان النهائي}\ifcorrection\ \textbf{-- التصحيح النموذجي}\fi}

\vspace{0.2cm}

{\Large\khatt\textbf{في الذكاء الاصطناعي}}

\end{center}

\rule{\textwidth}{1pt}

\ifcorrection\else
\textbf{الاسم واللقب:} \dotfill \quad \textbf{الفوج:} \makebox[3cm]{\dotfill}

\vspace{0.3cm}
\fi

\begin{center}

{\khatt\textbf{أسئلة متعددة الاختيارات (20 سؤالا -- 20 نقطة)}}

\end{center}

\rule{\textwidth}{1pt}

\textbf{التعليمة:} ضع علامة \textbf{$\times$} في المربع المقابل للإجابة الصحيحة. لكل سؤال إجابة صحيحة واحدة فقط.

\vspace{0.3cm}

\begin{enumerate}

\needspace{9\baselineskip}\item
\textbf{(الفصل 1) ما الذي يميّز الذكاء الاصطناعي «التوليدي» عن غيره من تطبيقات الذكاء الاصطناعي؟} \hfill \pts{1}
\begin{itemize}[leftmargin=1.6cm, labelsep=0.3cm, itemsep=2pt, topsep=4pt]
    \item[\emptybox\ \textbf{أ)}] أنه مخصص للحساب الرياضي ولا يتعامل مع اللغة.
    \item[\correctbox\ \textbf{ب)}] \correct{أنه ينتج محتوى جديدًا لم يكن موجودًا بهذا الشكل.}
    \item[\emptybox\ \textbf{ج)}] أنه لا يخطئ أبدًا لأنه يعتمد على قواعد ثابتة.
    \item[\emptybox\ \textbf{د)}] أنه يبحث فقط عن إجابة جاهزة مخزّنة مسبقًا.
    \item[\emptybox\ \textbf{هـ)}] أنه يعمل دون كهرباء أو اتصال بالإنترنت.
    \item[\emptybox\ \textbf{و)}] أنه يحفظ كل ما يكتبه المستخدم ويعيده كما هو.
    \item[\emptybox\ \textbf{ز)}] أنه يرتّب نتائج البحث حسب تاريخ نشرها.
    \item[\emptybox\ \textbf{ح)}] أنه جهاز عتادي يُركَّب داخل الهاتف لتسريعه.
\end{itemize}
\explanation{الذكاء الاصطناعي التوليدي نوع من الذكاء الاصطناعي ينتج محتوى جديدًا: نصًا أو صورة أو صوتًا أو جدولًا.}

%------------------------------------------------

\needspace{9\baselineskip}\item
\textbf{(الفصل 1) أجاب مساعد ذكاء اصطناعي بثقة أن بيتًا شعريًا لشاعر معيّن، ثم تبيّن أن الإجابة مختلَقة. ما اسم هذه الظاهرة؟} \hfill \pts{1}
\begin{itemize}[leftmargin=1.6cm, labelsep=0.3cm, itemsep=2pt, topsep=4pt]
    \item[\correctbox\ \textbf{أ)}] \correct{الهلوسة.}
    \item[\emptybox\ \textbf{ب)}] تعدّد الأنماط.
    \item[\emptybox\ \textbf{ج)}] التلخيص.
    \item[\emptybox\ \textbf{د)}] الانحياز الخوارزمي.
    \item[\emptybox\ \textbf{هـ)}] هندسة الأوامر.
    \item[\emptybox\ \textbf{و)}] الإفصاح.
    \item[\emptybox\ \textbf{ز)}] الانتحال العلمي.
    \item[\emptybox\ \textbf{ح)}] الشفافية.
\end{itemize}
\explanation{الهلوسة هي أن تقدّم الأداة معلومة غير صحيحة أو مختلَقة بطلاقة وثقة، ولهذا نتحقق دائمًا من مصدر موثوق.}

%------------------------------------------------

\needspace{9\baselineskip}\item
\textbf{(الفصل 1) أي أداة ذكاء اصطناعي توليدي تطوّرها شركة \textenglish{Anthropic}؟} \hfill \pts{1}
\begin{itemize}[leftmargin=1.6cm, labelsep=0.3cm, itemsep=2pt, topsep=4pt]
    \item[\emptybox\ \textbf{أ)}] \textenglish{Google}
    \item[\emptybox\ \textbf{ب)}] \textenglish{Perplexity}
    \item[\emptybox\ \textbf{ج)}] \textenglish{Gemini}
    \item[\emptybox\ \textbf{د)}] \textenglish{ChatGPT}
    \item[\emptybox\ \textbf{هـ)}] \textenglish{Alibaba}
    \item[\correctbox\ \textbf{و)}] \correct{\textenglish{Claude}}
    \item[\emptybox\ \textbf{ز)}] \textenglish{Qwen}
    \item[\emptybox\ \textbf{ح)}] \textenglish{DeepSeek}
\end{itemize}
\explanation{\textenglish{Claude} من \textenglish{Anthropic}، و\textenglish{ChatGPT} من \textenglish{OpenAI}، و\textenglish{Gemini} من \textenglish{Google}، و\textenglish{Qwen} من \textenglish{Alibaba}، و\textenglish{DeepSeek} من الشركة التي تحمل الاسم نفسه، و\textenglish{Perplexity} من \textenglish{Perplexity AI}.}

%------------------------------------------------

\needspace{9\baselineskip}\item
\textbf{(الفصل 1) ما الموقف الصحيح من دور الذكاء الاصطناعي في عمل أستاذ اللغة العربية؟} \hfill \pts{1}
\begin{itemize}[leftmargin=1.6cm, labelsep=0.3cm, itemsep=2pt, topsep=4pt]
    \item[\correctbox\ \textbf{أ)}] \correct{مساعد للأستاذ وليس بديلًا عنه، والأستاذ يبقى المسؤول.}
    \item[\emptybox\ \textbf{ب)}] وسيلة لتصحيح الامتحانات دون مراجعة الأستاذ.
    \item[\emptybox\ \textbf{ج)}] مسؤول عن الأخطاء بدلًا من الأستاذ.
    \item[\emptybox\ \textbf{د)}] أداة تغني عن المعاجم وكتب النحو المعتمدة.
    \item[\emptybox\ \textbf{هـ)}] بديل كامل عن الأستاذ في إعداد الدروس وتقديمها.
    \item[\emptybox\ \textbf{و)}] مصدر علمي معتمد لا تحتاج معلوماته إلى تحقق.
    \item[\emptybox\ \textbf{ز)}] أداة ممنوعة لا يجوز استعمالها في أي حال.
    \item[\emptybox\ \textbf{ح)}] أداة لا تفيد إلا أساتذة المواد العلمية.
\end{itemize}
\explanation{خلاصة الفصل الأول: الذكاء الاصطناعي مساعد للأستاذ وليس بديلًا عنه، والأستاذ يبقى المسؤول عن التحقق.}

%------------------------------------------------

\needspace{9\baselineskip}\item
\textbf{(الفصل 2) ما الترتيب الصحيح لمراحل تعلّم الآلة؟} \hfill \pts{1}
\begin{itemize}[leftmargin=1.6cm, labelsep=0.3cm, itemsep=2pt, topsep=4pt]
    \item[\emptybox\ \textbf{أ)}] النموذج $\leftarrow$ البيانات $\leftarrow$ التدريب $\leftarrow$ الاختبار $\leftarrow$ النتيجة
    \item[\emptybox\ \textbf{ب)}] النتيجة $\leftarrow$ البيانات $\leftarrow$ التدريب $\leftarrow$ النموذج $\leftarrow$ الاختبار
    \item[\emptybox\ \textbf{ج)}] التدريب $\leftarrow$ النموذج $\leftarrow$ البيانات $\leftarrow$ الاختبار $\leftarrow$ النتيجة
    \item[\emptybox\ \textbf{د)}] البيانات $\leftarrow$ النموذج $\leftarrow$ التدريب $\leftarrow$ النتيجة $\leftarrow$ الاختبار
    \item[\emptybox\ \textbf{هـ)}] التدريب $\leftarrow$ البيانات $\leftarrow$ النموذج $\leftarrow$ النتيجة $\leftarrow$ الاختبار
    \item[\emptybox\ \textbf{و)}] الاختبار $\leftarrow$ البيانات $\leftarrow$ التدريب $\leftarrow$ النموذج $\leftarrow$ النتيجة
    \item[\emptybox\ \textbf{ز)}] البيانات $\leftarrow$ الاختبار $\leftarrow$ التدريب $\leftarrow$ النموذج $\leftarrow$ النتيجة
    \item[\correctbox\ \textbf{ح)}] \correct{البيانات $\leftarrow$ التدريب $\leftarrow$ النموذج $\leftarrow$ الاختبار $\leftarrow$ النتيجة}
\end{itemize}
\explanation{تبدأ العملية بالبيانات، ثم التدريب الذي ينتج النموذج، ثم اختبار النموذج ببيانات جديدة، ثم النتيجة.}

%------------------------------------------------

\needspace{9\baselineskip}\item
\textbf{(الفصل 2) لماذا نختبر النموذج ببيانات جديدة لم يرها أثناء التدريب؟} \hfill \pts{1}
\begin{itemize}[leftmargin=1.6cm, labelsep=0.3cm, itemsep=2pt, topsep=4pt]
    \item[\emptybox\ \textbf{أ)}] لأن النموذج لا يقبل البيانات نفسها مرتين.
    \item[\emptybox\ \textbf{ب)}] لأن بيانات التدريب تُحذف بعد انتهاء التدريب.
    \item[\emptybox\ \textbf{ج)}] لأن الاختبار بالبيانات القديمة ممنوع قانونيًا.
    \item[\emptybox\ \textbf{د)}] لا سبب لذلك، فهي مجرد عادة عند المهندسين.
    \item[\emptybox\ \textbf{هـ)}] لزيادة حجم بيانات التدريب دون جهد إضافي.
    \item[\emptybox\ \textbf{و)}] لأن البيانات الجديدة أقصر فيكون الاختبار أسرع.
    \item[\correctbox\ \textbf{ز)}] \correct{لنعرف هل فهم القاعدة فعلًا، لا أنه يجيب من الذاكرة.}
    \item[\emptybox\ \textbf{ح)}] لتقليل استهلاك الكهرباء أثناء الاختبار.
\end{itemize}
\explanation{كما لا نختبر التلميذ بالجمل التي شُرحت له، نختبر النموذج ببيانات جديدة لنعرف هل فهم حقًا.}

%------------------------------------------------

\needspace{9\baselineskip}\item
\textbf{(الفصل 2) نموذج تدرّب على نصوص صحفية حديثة فقط، فعدّ كلمات فصيحة قديمة في الشعر الجاهلي أخطاءً. ما اسم هذه المشكلة؟} \hfill \pts{1}
\begin{itemize}[leftmargin=1.6cm, labelsep=0.3cm, itemsep=2pt, topsep=4pt]
    \item[\emptybox\ \textbf{أ)}] خطأ في هندسة الأمر لغياب الدور.
    \item[\emptybox\ \textbf{ب)}] خطأ في التنسيق الناتج عن غياب الجداول.
    \item[\correctbox\ \textbf{ج)}] \correct{الانحياز الخوارزمي الناتج عن بيانات غير ممثِّلة.}
    \item[\emptybox\ \textbf{د)}] خطأ في الأرقام عند قراءة الصور.
    \item[\emptybox\ \textbf{هـ)}] الانتحال الناتج عن نسخ النصوص الصحفية.
    \item[\emptybox\ \textbf{و)}] ضعف الإفصاح عن مصدر البيانات.
    \item[\emptybox\ \textbf{ز)}] نقص في تعدّد أنماط الأداة.
    \item[\emptybox\ \textbf{ح)}] الهلوسة الناتجة عن ضعف الاتصال بالإنترنت.
\end{itemize}
\explanation{النموذج لم يرَ هذا النوع من النصوص أثناء التدريب، فبياناته غير ممثِّلة، فجاءت نتائجه منحازة وغير دقيقة.}

%------------------------------------------------

\needspace{9\baselineskip}\item
\textbf{(الفصل 3) في الأمر: «تصرّف كأستاذ لغة عربية...»، ما العنصر الذي تمثّله هذه العبارة؟} \hfill \pts{1}
\begin{itemize}[leftmargin=1.6cm, labelsep=0.3cm, itemsep=2pt, topsep=4pt]
    \item[\emptybox\ \textbf{أ)}] القيود
    \item[\emptybox\ \textbf{ب)}] الإفصاح
    \item[\emptybox\ \textbf{ج)}] الكلمات المفتاحية
    \item[\correctbox\ \textbf{د)}] \correct{الدور}
    \item[\emptybox\ \textbf{هـ)}] السياق
    \item[\emptybox\ \textbf{و)}] النمط
    \item[\emptybox\ \textbf{ز)}] التنسيق
    \item[\emptybox\ \textbf{ح)}] المهمة
\end{itemize}
\explanation{الدور يحدد من تكون الأداة وأي خبرة تعتمد عليها.}

%------------------------------------------------

\needspace{9\baselineskip}\item
\textbf{(الفصل 3) أضاف أستاذ إلى أمره: «في 150 كلمة، بلغة بسيطة». ما العنصر الذي أضافه؟} \hfill \pts{1}
\begin{itemize}[leftmargin=1.6cm, labelsep=0.3cm, itemsep=2pt, topsep=4pt]
    \item[\correctbox\ \textbf{أ)}] \correct{القيود}
    \item[\emptybox\ \textbf{ب)}] الشفافية
    \item[\emptybox\ \textbf{ج)}] البيانات
    \item[\emptybox\ \textbf{د)}] الدور
    \item[\emptybox\ \textbf{هـ)}] النموذج
    \item[\emptybox\ \textbf{و)}] السياق
    \item[\emptybox\ \textbf{ز)}] المهمة
    \item[\emptybox\ \textbf{ح)}] التنسيق
\end{itemize}
\explanation{القيود هي الشروط التي يجب احترامها في الإجابة، مثل عدد الكلمات وبساطة اللغة.}

%------------------------------------------------

\needspace{9\baselineskip}\item
\textbf{(الفصل 3) أي أمر يحتوي على العناصر الخمسة للأمر الجيد؟} \hfill \pts{1}
\begin{itemize}[leftmargin=1.6cm, labelsep=0.3cm, itemsep=2pt, topsep=4pt]
    \item[\emptybox\ \textbf{أ)}] اشرح الحال.
    \item[\correctbox\ \textbf{ب)}] \correct{تصرّف كأستاذ نحو. اشرح الحال لتلاميذ الثانية متوسط، في 150 كلمة، في شكل تعريف ثم أمثلة ثم خلاصة.}
    \item[\emptybox\ \textbf{ج)}] اشرح الحال في جدول.
    \item[\emptybox\ \textbf{د)}] اشرح الحال وترجمه واكتب اختبارًا وصحّحه.
    \item[\emptybox\ \textbf{هـ)}] أنا أستاذ. اشرح الحال.
    \item[\emptybox\ \textbf{و)}] اشرح الحال في 150 كلمة.
    \item[\emptybox\ \textbf{ز)}] اكتب شيئًا جميلًا عن النحو.
    \item[\emptybox\ \textbf{ح)}] تصرّف كأستاذ نحو. اشرح الحال.
\end{itemize}
\explanation{الأمر الأول يجمع الدور والمهمة والسياق والقيود والتنسيق. الأوامر الأخرى ينقصها عنصر أو أكثر، أو تخلط مهام كثيرة.}

%------------------------------------------------

\needspace{9\baselineskip}\item
\textbf{(الفصل 3) لم تكن أستاذة راضية عن أول إجابة من الأداة. ما التصرف الصحيح؟} \hfill \pts{1}
\begin{itemize}[leftmargin=1.6cm, labelsep=0.3cm, itemsep=2pt, topsep=4pt]
    \item[\emptybox\ \textbf{أ)}] كتابة الأمر بالأحرف الكبيرة.
    \item[\emptybox\ \textbf{ب)}] إضافة مهام أخرى كثيرة إلى الأمر نفسه.
    \item[\emptybox\ \textbf{ج)}] قبول الإجابة الأولى لأنها الأفضل دائمًا.
    \item[\emptybox\ \textbf{د)}] حذف السياق والقيود لتقصير الأمر.
    \item[\correctbox\ \textbf{هـ)}] \correct{تعديل الأمر وإضافة ما ينقصه، ثم إعادة المحاولة.}
    \item[\emptybox\ \textbf{و)}] التوقف نهائيًا عن استعمال الأداة.
    \item[\emptybox\ \textbf{ز)}] تكرار الأمر نفسه حرفيًا عشر مرات.
    \item[\emptybox\ \textbf{ح)}] نسخ إجابة أداة أخرى دون قراءتها.
\end{itemize}
\explanation{قبول أول إجابة من الأخطاء الشائعة. نبني الأمر تدريجيًا ونحسّنه بعد قراءة الإجابة.}

%------------------------------------------------

\needspace{9\baselineskip}\item
\textbf{(الفصل 4) قدّم أستاذ معلومة خاطئة أخذها من أداة دون تحقق. من المسؤول؟} \hfill \pts{1}
\begin{itemize}[leftmargin=1.6cm, labelsep=0.3cm, itemsep=2pt, topsep=4pt]
    \item[\emptybox\ \textbf{أ)}] الأداة وحدها لأنها ولّدت المعلومة.
    \item[\correctbox\ \textbf{ب)}] \correct{الأستاذ، لأن الاستعانة بالأداة لا تنقل المسؤولية.}
    \item[\emptybox\ \textbf{ج)}] مدير المؤسسة لأنه لم يمنع الأداة.
    \item[\emptybox\ \textbf{د)}] التلاميذ لأنهم لم ينتبهوا للخطأ.
    \item[\emptybox\ \textbf{هـ)}] الشركة المطوّرة للأداة وحدها.
    \item[\emptybox\ \textbf{و)}] لا أحد، لأن الأخطاء الرقمية لا تُحاسَب.
    \item[\emptybox\ \textbf{ز)}] المسؤولية مقسومة بالتساوي بين الأستاذ والأداة.
    \item[\emptybox\ \textbf{ح)}] محرك البحث الذي أوصل إلى الأداة.
\end{itemize}
\explanation{مبدأ المسؤولية: المستخدم مسؤول عن كل معلومة يعتمد عليها، والأداة لا تتحمل المسؤولية.}

%------------------------------------------------

\needspace{9\baselineskip}\item
\textbf{(الفصل 4) أنشأ طالب مسودة بالأداة، ثم راجعها وتحقق منها وصرّح بذلك. كيف يُصنَّف هذا الاستعمال؟} \hfill \pts{1}
\begin{itemize}[leftmargin=1.6cm, labelsep=0.3cm, itemsep=2pt, topsep=4pt]
    \item[\emptybox\ \textbf{أ)}] مرفوض دائمًا.
    \item[\emptybox\ \textbf{ب)}] انحياز خوارزمي.
    \item[\correctbox\ \textbf{ج)}] \correct{مقبول بشروط.}
    \item[\emptybox\ \textbf{د)}] هلوسة.
    \item[\emptybox\ \textbf{هـ)}] غير مسؤول.
    \item[\emptybox\ \textbf{و)}] مقبول دون أي شرط.
    \item[\emptybox\ \textbf{ز)}] ممنوع في كل المؤسسات.
    \item[\emptybox\ \textbf{ح)}] انتحال مرفوض.
\end{itemize}
\explanation{إنشاء مسودة بالأداة ثم مراجعتها والتحقق منها والإفصاح استعمال «مقبول بشروط» في مقياس الدرس.}

%------------------------------------------------

\needspace{9\baselineskip}\item
\textbf{(الفصل 4) نسخ طالب بحثًا كاملًا من أداة وقدّمه باسمه دون تصريح. ما اسم هذا السلوك؟} \hfill \pts{1}
\begin{itemize}[leftmargin=1.6cm, labelsep=0.3cm, itemsep=2pt, topsep=4pt]
    \item[\emptybox\ \textbf{أ)}] الشفافية.
    \item[\emptybox\ \textbf{ب)}] الاستعمال المقبول.
    \item[\correctbox\ \textbf{ج)}] \correct{الانتحال بالذكاء الاصطناعي.}
    \item[\emptybox\ \textbf{د)}] التلخيص.
    \item[\emptybox\ \textbf{هـ)}] التعلّم الآلي.
    \item[\emptybox\ \textbf{و)}] تقييم المصادر.
    \item[\emptybox\ \textbf{ز)}] هندسة الأوامر.
    \item[\emptybox\ \textbf{ح)}] الإفصاح.
\end{itemize}
\explanation{الانتحال بالذكاء الاصطناعي: أن ينسخ الطالب نصًا أنتجته الأداة ويقدّمه على أنه عمله الشخصي.}

%------------------------------------------------

\needspace{9\baselineskip}\item
\textbf{(الفصل 5) أي مما يلي سؤال بحث محدد وليس موضوعًا عامًا؟} \hfill \pts{1}
\begin{itemize}[leftmargin=1.6cm, labelsep=0.3cm, itemsep=2pt, topsep=4pt]
    \item[\emptybox\ \textbf{أ)}] اللغة العربية في المدرسة.
    \item[\emptybox\ \textbf{ب)}] الأدب الجاهلي.
    \item[\emptybox\ \textbf{ج)}] تعليم النحو.
    \item[\correctbox\ \textbf{د)}] \correct{ما الصعوبات التي يواجهها تلاميذ المتوسط في الإعراب؟}
    \item[\emptybox\ \textbf{هـ)}] البلاغة العربية.
    \item[\emptybox\ \textbf{و)}] الذكاء الاصطناعي.
    \item[\emptybox\ \textbf{ز)}] الذكاء الاصطناعي والتعليم.
    \item[\emptybox\ \textbf{ح)}] التكنولوجيا الحديثة.
\end{itemize}
\explanation{موضوع البحث مجال عام وواسع، أما سؤال البحث فسؤال محدد نريد الإجابة عنه.}

%------------------------------------------------

\needspace{9\baselineskip}\item
\textbf{(الفصل 5) أي مصدر هو الأكثر موثوقية حسب سلّم الدرس؟} \hfill \pts{1}
\begin{itemize}[leftmargin=1.6cm, labelsep=0.3cm, itemsep=2pt, topsep=4pt]
    \item[\emptybox\ \textbf{أ)}] موقع مجهول لا يُعرف كاتبه.
    \item[\emptybox\ \textbf{ب)}] منشور على شبكة اجتماعية.
    \item[\emptybox\ \textbf{ج)}] أول نتيجة في محرك البحث.
    \item[\emptybox\ \textbf{د)}] مقال صحفي.
    \item[\emptybox\ \textbf{هـ)}] محتوى أنتجه الذكاء الاصطناعي.
    \item[\emptybox\ \textbf{و)}] إعلان لدورة تكوينية مدفوعة.
    \item[\emptybox\ \textbf{ز)}] مدونة شخصية.
    \item[\correctbox\ \textbf{ح)}] \correct{كتاب محقق أو مقال علمي محكّم.}
\end{itemize}
\explanation{المصادر الأكاديمية في أعلى السلّم، ثم مواقع المؤسسات الرسمية، ثم المقال الصحفي، ثم المدونة، وفي الأسفل محتوى الذكاء الاصطناعي.}

%------------------------------------------------

\needspace{9\baselineskip}\item
\textbf{(الفصل 5) صفحة تمدح طريقة لتعليم النحو وتنتهي برابط لشراء دورة. أي معيار يكشف ضعفها؟} \hfill \pts{1}
\begin{itemize}[leftmargin=1.6cm, labelsep=0.3cm, itemsep=2pt, topsep=4pt]
    \item[\emptybox\ \textbf{أ)}] الكاتب
    \item[\emptybox\ \textbf{ب)}] النمط
    \item[\emptybox\ \textbf{ج)}] التنسيق
    \item[\emptybox\ \textbf{د)}] الدور
    \item[\emptybox\ \textbf{هـ)}] الدقة
    \item[\correctbox\ \textbf{و)}] \correct{الهدف}
    \item[\emptybox\ \textbf{ز)}] الجهة
    \item[\emptybox\ \textbf{ح)}] التاريخ
\end{itemize}
\explanation{معيار الهدف يسأل: هل هدف المصدر الإعلام، أم الإقناع، أم البيع؟}

%------------------------------------------------

\needspace{9\baselineskip}\item
\textbf{(الفصل 6) ما المقصود بالأداة «متعددة الأنماط»؟} \hfill \pts{1}
\begin{itemize}[leftmargin=1.6cm, labelsep=0.3cm, itemsep=2pt, topsep=4pt]
    \item[\emptybox\ \textbf{أ)}] أداة تعمل على الهاتف والحاسوب معًا.
    \item[\emptybox\ \textbf{ب)}] أداة تتصل بعدة محركات بحث في الوقت نفسه.
    \item[\emptybox\ \textbf{ج)}] أداة يستعملها عدة أشخاص في الحساب نفسه.
    \item[\emptybox\ \textbf{د)}] أداة تعمل بعدة لغات لكن بالنص فقط.
    \item[\emptybox\ \textbf{هـ)}] أداة لها عدة أسماء تجارية مختلفة.
    \item[\emptybox\ \textbf{و)}] أداة تُحدَّث عدة مرات في السنة.
    \item[\correctbox\ \textbf{ز)}] \correct{أداة تفهم أكثر من نمط وتنتج أكثر من نمط.}
    \item[\emptybox\ \textbf{ح)}] أداة تقدّم عدة إجابات نصية للسؤال نفسه.
\end{itemize}
\explanation{النمط شكل المعلومة (نص، صورة، صوت، فيديو). الأداة متعددة الأنماط تفهم أكثر من نمط وتنتج أكثر من نمط.}

%------------------------------------------------

\needspace{9\baselineskip}\item
\textbf{(الفصل 6) صوّر أستاذ صفحة مخطوطة وطلب كتابتها نصًا رقميًا. ما نوع هذا التحويل؟} \hfill \pts{1}
\begin{itemize}[leftmargin=1.6cm, labelsep=0.3cm, itemsep=2pt, topsep=4pt]
    \item[\emptybox\ \textbf{أ)}] صوت $\leftarrow$ فيديو
    \item[\emptybox\ \textbf{ب)}] نص $\leftarrow$ صوت
    \item[\emptybox\ \textbf{ج)}] نص $\leftarrow$ صورة
    \item[\correctbox\ \textbf{د)}] \correct{صورة $\leftarrow$ نص}
    \item[\emptybox\ \textbf{هـ)}] مخطط $\leftarrow$ صوت
    \item[\emptybox\ \textbf{و)}] نص $\leftarrow$ مخطط
    \item[\emptybox\ \textbf{ز)}] فيديو $\leftarrow$ صورة
    \item[\emptybox\ \textbf{ح)}] صوت $\leftarrow$ نص
\end{itemize}
\explanation{المدخل صورة (صفحة المخطوطة) والمخرج نص رقمي: «صورة $\leftarrow$ نص».}

%------------------------------------------------

\needspace{9\baselineskip}\item
\textbf{(الفصل 6) أي مجال تخطئ فيه الأداة غالبًا عند قراءة صورة أو مخطط بياني؟} \hfill \pts{1}
\begin{itemize}[leftmargin=1.6cm, labelsep=0.3cm, itemsep=2pt, topsep=4pt]
    \item[\emptybox\ \textbf{أ)}] اسم الجهاز الذي التُقطت به الصورة.
    \item[\emptybox\ \textbf{ب)}] لا تخطئ أبدًا في قراءة الصور.
    \item[\emptybox\ \textbf{ج)}] عدد الألوان المستعملة في الصورة.
    \item[\emptybox\ \textbf{د)}] تاريخ التقاط الصورة بالهاتف.
    \item[\correctbox\ \textbf{هـ)}] \correct{الأرقام: قد تقرأ رقمًا خطأ أو تخلط بين عمودين.}
    \item[\emptybox\ \textbf{و)}] حجم ملف الصورة على الحاسوب.
    \item[\emptybox\ \textbf{ز)}] اتجاه الصورة أفقيًا أو عموديًا فقط.
    \item[\emptybox\ \textbf{ح)}] اسم الملف الذي حُفظت به الصورة.
\end{itemize}
\explanation{تخطئ الأداة غالبًا في الأرقام والعلاقات والخط العربي، لذلك نقارن إجابتها دائمًا بالصورة الأصلية.}

\end{enumerate}

\vfill
\begin{center}
{\khatt بالتوفيق}
\end{center}

\end{document}
